\infographicchapter{Interventional Radiology}{interventional_radiology.png}
\label{ch:infographic-interventional-radiology}

\section{Overview}

Interventional radiology (IR) is a medical specialty that diagnoses and treats diseases throughout the body using tiny, pinhole-sized skin openings instead of large surgical cuts. Rather than relying on one single scanner, interventional radiologists use live imaging---most often X-ray fluoroscopy, CT, ultrasound, or MRI---to guide hair-thin wires, slim tubes (catheters), and specialized needles straight to deep organs and blood vessels with pinpoint accuracy.

IR treats a huge range of conditions across nearly every body system. Common procedures include opening narrowed or blocked blood vessels with balloons and stents, blocking abnormal bleeding or shrinking uterine fibroids (embolization), destroying liver, kidney, lung, or bone tumors with targeted heat, freezing, or tiny beads, taking precise needle biopsies, draining deep fluid pockets or blocked bile and kidney tubes, placing long-term IV ports, and stabilizing painful spine fractures.

Compared with traditional open surgery, IR typically causes less pain, smaller scars, lower complication risks, shorter hospital stays, and faster recovery---often needing only local numbing medicine and relaxing sedation instead of full general anesthesia. Its limits include needing specialized imaging rooms and expert teams, and it cannot replace open surgery when large organs need reconstruction or wide removal. Safety focuses on keeping X-ray doses as low as possible (ALARA), checking contrast allergies and kidney health, managing blood thinners to prevent bleeding, and using strict operating-room sterile technique to prevent infection.

\section{Underlying Physical Principles}

\begin{itemize}
    \item \textbf{Choosing the Best Live Guidance:} Interventional radiologists pick the imaging tool that best fits the job---live pulsed X-rays (fluoroscopy) for blood vessels, ultrasound for soft organs near the surface, CT for deep bone or lung targets, or MRI when radiation-free soft-tissue detail is needed.
    \item \textbf{The Seldinger Technique (Needle-to-Catheter Swap):} Building on Sven Ivar Seldinger's 1953 breakthrough and Charles Dotter's 1964 catheter angioplasty, IR safely trades a thin needle for working tools without surgical cuts: the doctor enters the vessel or cavity with a hollow needle, threads a soft flexible guidewire inside, removes the needle, and slides a protective sheath and catheter smoothly over the wire.
    \item \textbf{Erasing Background Bone to See Vessels Clearly:} During blood-vessel procedures, digital subtraction angiography (DSA) takes a quick baseline X-ray before dye is injected and digitally subtracts bones and background organs from the live frames. This leaves crystal-clear maps of flowing blood vessels, and ``roadmapping'' overlays the live moving wire right on top of that vessel map.
    \item \textbf{3D Live Maps in the Procedure Room:} Modern C-arm X-ray machines can spin around you in seconds to create a 3D cone-beam CT scan right on the procedure table, blending it with earlier CT or MRI scans so the doctor can track needles and catheters from every angle.
    \item \textbf{Pinpoint Treatment Delivery:} Once the needle or catheter reaches the exact target, the doctor delivers focused treatment---stretching a tight vessel with a tiny balloon or stent, blocking a bleeding artery with microscopic beads or coils, or destroying a tumor with focused heat (radiofrequency or microwave) or extreme cold (cryoablation).
\end{itemize}

\section{Imaging Mechanism}

\begin{enumerate}
    \item \textbf{Room Setup and Sterile Prep:} You lie on a specialized angiography or CT table while the skin over the tiny entry spot is cleaned thoroughly and covered with sterile drapes. The C-arm X-ray detector, sterile-covered ultrasound probe, or CT scanner is lined up and optionally matched to your earlier CT or MRI scans.
    \item \textbf{Timed Contrast and Live Tracking:} A precision power injector delivers timed contrast dye (or harmless carbon dioxide gas when iodine must be avoided) synchronized with live X-ray video, 3D rotational C-arm sweeps, or real-time CT and ultrasound needle tracking.
    \item \textbf{Instant Digital Capture:} Flat-panel X-ray arrays, CT detectors, or ultrasound probes capture live signals and stream them instantly into high-speed digital computers.
    \item \textbf{Building Subtracted and 3D Maps:} Workstations subtract background bone, correct for tiny movements, rebuild 3D cone-beam CT slices, and fuse vessel or tumor maps together in seconds.
    \item \textbf{Multi-Screen Live Navigation:} Large overhead screens display live X-ray video, vessel roadmaps, 3D slices, ultrasound, and your heart and blood-pressure monitors all at once, guiding every millimeter of the procedure.
\end{enumerate}

\section{Key Technical Parameters}

\noindent\begin{tabularx}{\textwidth}{@{}l X@{}}
\toprule
\textbf{Parameter} & \textbf{Typical Value / Description} \\
\midrule
Physical Principle & Pinhole-entry diagnosis and treatment guided live by X-ray fluoroscopy/DSA, CT, ultrasound, or MRI \\
Energy / Signal Source & Multi-modality guidance: pulsed X-ray angiography, CT X-rays, 2--15 MHz ultrasound, or 1.5--3.0 T MRI \\
Typical Frequency / Energy & 60--125 kVp (X-ray/CT guidance); 2--15 MHz (ultrasound); 1.5--3.0 T (interventional MRI) \\
Spatial Resolution & Very sharp 0.15--0.5 mm detail on flat-panel angiography and ultrasound; 0.5--1.0 mm on CT guidance \\
Typical Scan Duration & 30--60 minutes for biopsies, drainages, or IV ports; 1--3 hours for complex vessel or tumor treatments \\
Contrast Agent & IV iodine contrast (or harmless carbon dioxide gas if you have kidney failure or severe allergy), tube contrast, or ultrasound microbubbles \\
Ionizing Radiation & Depends on the guide used: zero with ultrasound or MRI; moderate to higher localized doses with long X-ray or CT procedures, closely tracked under ALARA \\
\bottomrule
\end{tabularx}

\section{Common Symptoms That Prompt This Scan}
\begin{itemize}
    \item \textbf{Active Bleeding or Blocked Blood Vessels:} Severe leg pain when walking, sudden limb numbness or coldness, acute stroke symptoms, or internal bleeding that requires pinpoint pinhole treatment.
    \item \textbf{A Deep Lump or Fluid Collection Needing Diagnosis or Drainage:} Fever and localized pain from a deep abscess, yellowing skin (jaundice) from blocked bile ducts, flank pain from a blocked kidney, or a suspicious mass needing a needle biopsy.
    \item \textbf{Heavy Menstrual Bleeding or Pelvic Pressure:} Painful, heavy periods or bladder pressure caused by uterine fibroids eligible for minimally invasive embolization.
    \item \textbf{Severe Pain From a Collapsed Spine Bone or Localized Tumor:} Debilitating mid- or lower-back pain from a spinal fracture, or a liver, kidney, or lung tumor scheduled for targeted heat or freezing ablation.
\end{itemize}

\section{Clinical Usages}

\subsection{Primary Usages}

\begin{itemize}
    \item \textbf{Opening or Sealing Blood Vessels:} Stretches narrowed arteries, places stents, and pulls out blood clots for leg-artery disease or acute stroke; seals off bleeding vessels after injury, treats tangled vessel malformations, and shrinks symptomatic uterine fibroids (benign uterine growths present in about 70\% of women and symptomatic in roughly 25\% of reproductive-age women) via embolization.
    \item \textbf{Minimally Invasive Cancer Care:} Destroys liver, kidney, lung, and bone tumors using targeted heat or freezing needles (ablation), and delivers tiny chemotherapy or radiation beads directly into liver tumor arteries (TACE and TARE) while sparing healthy tissue.
    \item \textbf{Guided Biopsies and Fluid Drains:} Collects tiny, accurate tissue cores from deep organs or lymph nodes and places slim drain tubes into infected fluid pockets (abscesses) in the belly, chest, or soft tissues.
    \item \textbf{Relieving Blocked Bile and Kidney Tubes:} Places gentle drain tubes or stents through the skin to relieve blocked bile ducts (jaundice) or blocked kidney drainage tubes (nephrostomy).
    \item \textbf{Long-Term IV Access and Clot Treatment:} Places chemotherapy ports, dialysis lines, and clot-catching filters (IVC filters), and dissolves or removes dangerous deep-vein or lung blood clots.
\end{itemize}

\subsection{Secondary Usages}

\begin{itemize}
    \item \textbf{Treating High Liver Vein Pressure:} Creates an internal liver bypass channel (TIPS) or seals swollen stomach veins (BRTO) to stop severe liver-related bleeding or stubborn belly fluid buildup.
    \item \textbf{Bone and Spine Relief:} Stabilizes painful collapsed spine bones with medical bone cement (vertebroplasty or kyphoplasty) and treats painful benign bone spots (osteoid osteomas) with targeted heat.
    \item \textbf{Feeding Tubes and Lymph Leaks:} Places stomach feeding tubes safely under imaging guidance and seals rare chest lymph-fluid leaks.
    \item \textbf{Pelvic, Fertility, and Vein Care:} Opens blocked fallopian tubes non-surgically to treat infertility (fallopian tube recanalization), eases chronic pelvic pain from swollen internal veins (pelvic congestion syndrome), and treats testicular varicoceles or venous malformations with embolization or vein-hardening injections (sclerotherapy).
    \item \textbf{Newer Targeted Therapies:} Redirects liver blood flow before liver surgery so healthy liver grows larger, shrinks enlarged prostates (prostatic artery embolization), and eases chronic knee osteoarthritis pain (genicular artery embolization).
    \item \textbf{Aneurysm Stent-Grafts, Vein Sealing, and Frail or Pediatric Care:} Reinforces dangerously bulging aortic aneurysms from inside the vessel with fabric-lined stent-grafts (EVAR), seals painful varicose veins with thin laser fibers, drains infected gallbladders (cholecystostomy) when surgery is too risky, and provides gentle venous access and needle biopsies---the most common IR procedures in children.
\end{itemize}

\section{Advantages}

\begin{itemize}
    \item \textbf{Tiny Pinhole Entry:} Reaches deep blood vessels and organs through a tiny nick in the skin rather than a large surgical cut, dramatically reducing tissue trauma.
    \item \textbf{Less Pain and Quicker Recovery:} Typically leads to less bleeding, less postoperative pain, fewer wound problems, shorter hospital stays, and a faster return to daily life than open surgery.
    \item \textbf{Often Avoids General Anesthesia:} Many procedures need only local numbing medicine and relaxing sedation, making treatment safer for older adults or people with heart and lung conditions.
    \item \textbf{Pinpoint Live Accuracy:} Real-time X-ray, ultrasound, CT, or MRI lets the doctor see exact targets and steer safely around nearby vital organs.
    \item \textbf{Handles Emergencies and Everyday Care:} Quickly stops life-threatening bleeding, opens blocked vessels, drains infections, takes biopsies, and treats tumors across nearly every organ system.
\end{itemize}

\section{Disadvantages}

\begin{itemize}
    \item \textbf{Small Procedural Risks:} Any internal puncture carries site-specific risks such as bleeding, bruising or a small vessel bulge (pseudoaneurysm) at the entry site, vessel wall injury, beads drifting to an unintended spot, infection, a small lung air leak (pneumothorax) during chest procedures, or accidental injury to nearby organs.
    \item \textbf{X-Ray and Contrast Exposure:} X-ray- and CT-guided procedures involve ionizing radiation and often iodine contrast dye, carrying small risks of allergic reactions or temporary kidney strain.
    \item \textbf{Cannot Replace Every Surgery:} Open or laparoscopic surgery is still required when large tumors must be removed in one piece, complex organs must be rebuilt, or no safe straight path exists for a needle.
    \item \textbf{Requires Specialized Centers:} Relies on advanced angiography rooms, specialized tools, and highly trained subspecialty teams, which may not be available at smaller clinics.
    \item \textbf{Not Safe in Every Situation:} Severe uncontrolled bleeding disorders, active skin infection right over the planned entry spot, severe untreated contrast allergy, or inability to lie flat safely can delay or prevent a procedure.
\end{itemize}

\section{Contraindications}

\subsection{Absolute}

\begin{itemize}
    \item \textbf{No Safe Path or Uncontrolled Bleeding:} A procedure cannot be done if vital organs or major blood vessels block every safe needle or catheter path, if severe uncontrolled clotting problems prevent stopping bleeding safely during a high-risk intervention, or if the patient declines consent.
\end{itemize}

\subsection{Relative}

\begin{itemize}
    \item \textbf{Treatable Clotting or Skin Issues:} Mild-to-moderate clotting problems or low platelets that can be safely corrected with blood products or pausing blood thinners, or a localized skin infection at the entry site (where a different clean entry spot can be chosen).
    \item \textbf{Contrast, Pregnancy, and Positioning Concerns:} Past severe iodine contrast allergy (often managed with allergy premedication or harmless carbon dioxide gas), advanced kidney disease, pregnancy during X-ray- or CT-guided procedures, or trouble lying flat without anesthesia help.
\end{itemize}

\section{Patient Safety Guidelines}

\begin{itemize}
    \item \textbf{Safety Time-Out and Blood-Thinner Check:} Before starting, the entire team performs a formal ``time-out'' to verify your identity, procedure, and exact side, while confirming kidney labs, allergy history, clotting numbers, and blood-thinner pauses.
    \item \textbf{Operating-Room Sterile Care:} Staff use full sterile gowns, drapes, skin antiseptics, sterile ultrasound probe covers, and targeted preventive antibiotics when needed to keep infection risk as low as possible.
    \item \textbf{Radiation Protection and Dose Tracking:} Teams strictly follow ALARA X-ray safety---using pulsed beams, tight focus, and varied camera angles---and schedule a skin checkup if a long procedure exceeds high dose thresholds (such as 5 Gy).
    \item \textbf{Pregnancy and Contrast Safety:} Patients who could be pregnant are screened before X-ray or CT guidance. Teams use the smallest contrast volume needed, keep you well hydrated, and keep emergency allergy medicines ready.
    \item \textbf{Safe Closure and Recovery:} At the end of a vessel procedure, staff seal the tiny entry hole with firm pressure or a tiny closure device, then watch your vital signs, pulses, and puncture site closely as sedation wears off.
\end{itemize}

\section{Prerequisite Tests}

\begin{itemize}
    \item \textbf{Clotting and Blood Count Labs:} Depending on your procedure's bleeding risk, your doctor may check clotting numbers (INR/PT), platelets, and red blood cell counts---routinely for higher-risk solid-organ biopsies, TIPS, or bile drains, and selectively for simpler procedures.
    \item \textbf{Kidney Function Blood Test:} A recent creatinine and eGFR blood test is checked before procedures that use IV iodine contrast dye.
    \item \textbf{Recent Diagnostic Scans:} Your interventional radiologist reviews recent CT, MRI, or ultrasound scans to map the exact target, plan the safest needle or catheter path, and pick the right device sizes.
    \item \textbf{Health, Allergy, and Pregnancy Review:} A quick pre-procedure check covers contrast allergies, airway fitness for sedation, signs of active infection, and a pregnancy test when applicable before X-ray or CT guidance.
\end{itemize}

\section{Preparing for the Scan}

\begin{itemize}
    \item \textbf{Fasting and Hydration:} If you will receive sedation or anesthesia, do not eat solid food for 6--8 hours and stop clear liquids 2 hours beforehand. Drink plenty of water earlier in the day before contrast procedures if instructed.
    \item \textbf{Blood Thinners and Daily Medicines:} Follow your IR clinic's exact plan for pausing or bridging blood thinners and aspirin. Adjust diabetes pills or insulin for fasting, and take approved morning pills (like blood-pressure medicine) with just a tiny sip of water.
    \item \textbf{Clean Skin and Simple Clothing:} Shower beforehand with antibacterial soap if asked, skip lotions or powders near the planned entry area, remove jewelry and metal accessories, and change into a hospital gown.
    \item \textbf{Driver and Early Arrival:} Arrange a responsible adult to drive you home and stay with you if you have outpatient sedation. Bring prior scan records and your medicine/allergy list, and arrive early for an IV line, site marking, consent, and safety time-out.
\end{itemize}

\section{During Scan}

\begin{itemize}
    \item \textbf{Getting Settled and Relaxed:} You will lie flat on a cushioned angiography or CT table while the entry skin---often at the wrist or groin for blood-vessel access, or directly over the target organ---is cleaned and covered with sterile drapes. Most procedures use local numbing medicine plus relaxing IV sedation so you feel drowsy and comfortable.
    \item \textbf{Staying Still and Brief Breath-Holds:} Keep your body still---especially keeping hands and arms away from the sterile drapes---and follow short breath-hold cues during X-ray or CT sweeps so pictures stay razor-sharp.
    \item \textbf{What You Will Feel and Hear:} Expect a brief stinging pinch when local numbing medicine goes in, followed mostly by dull pressure or gentle tugging rather than sharp pain. If iodine contrast dye is injected into a vessel, a warm flush and metallic taste lasting 15--30 seconds are completely normal, along with quiet machine hums.
    \item \textbf{Constant Nurse Care:} A dedicated nurse watches your heart rate, blood pressure, and oxygen continuously. Speak up right away if you feel sharp pain, itching, shortness of breath, chest discomfort, nausea, or numbness in an arm or leg.
    \item \textbf{How Long It Takes:} Routine biopsies, fluid drainages, or IV ports usually take 30--60 minutes. More complex vessel embolizations or tumor treatments typically take 1--3 hours, followed by monitored recovery.
\end{itemize}

\section{Normal Appearance}

\begin{itemize}
    \item \textbf{Healthy Blood-Vessel Maps (DSA):} Arteries taper smoothly with brisk, even contrast flow, a smooth brief tissue blush, and timely vein return---with no leaking dye, wall tears, or abnormal early vein shortcuts.
    \item \textbf{Clean Needle and Drain Placement:} Needles and coiled drain tubes line up directly inside the target lump or fluid pocket, staying safely clear of lungs, bowel, and major blood vessels.
    \item \textbf{Successful Treatment Signs:} Widened vessels flow briskly after a balloon or stent, bleeding or tumor-feeding arteries stop flowing cleanly after embolization with no stray drift, and a tumor ablation zone evenly covers the tumor plus a healthy 5--10 mm safety rim.
\end{itemize}

\section{Observing for Abnormal Patterns}

\begin{itemize}
    \item \textbf{Active Bleeding and Vessel Bulges:} On live vessel X-rays, a cloudy puddle of contrast that keeps growing confirms active bleeding, while a round contrast-filled pouch hanging off an artery signals a pseudoaneurysm that needs sealing.
    \item \textbf{Narrowed, Blocked, or Torn Vessels:} Pinched vessel segments, abrupt cutoffs from blood clots, slow flow, and extra detour vessels point to significant blockages, while a thin flap splitting the flow channel indicates a vessel-wall tear (dissection).
    \item \textbf{Abnormal Vessel Shortcuts and Tumor Supply:} Veins lighting up too early during the artery phase reveal an abnormal shortcut (fistula or malformation)---crucial to spot before liver tumor beads are given so beads do not drift to the lungs---while a bright, lingering stain highlights blood-rich tumors.
    \item \textbf{Stray Bead Backflow and Remaining Tumor Edge:} Contrast drifting backward during embolization warns the doctor to stop so beads do not enter healthy vessels. After tumor ablation, any remaining bright rim on the edge signals leftover tumor needing extra touch-up.
    \item \textbf{Harmless Look-Alikes:} Breathing, bowel movement, or heartbeats between baseline and dye frames create dark-and-white shadow streaks that can mimic or hide bleeding. Gentle rippled ``standing waves'' or temporary vessel spasm around a catheter tip can also look like fixed narrowing until the catheter relaxes or medicine opens the vessel.
\end{itemize}

\section{Dealing with Uncertainties}

\begin{itemize}
    \item \textbf{Movement Streaks on Vessel Maps:} Breathing or bowel motion can blur subtracted vessel pictures and hide small leaks. Checking the unsubtracted raw frames, digitally shifting the image back into alignment, and coaching brief breath-holds restore clarity.
    \item \textbf{Temporary Spasm and On-and-Off Bleeding:} Temporary artery spasm can mimic permanent narrowing, and digestive bleeding that pauses during the test can look normal even when bleeding has been happening.
    \item \textbf{Unclear Biopsy Samples:} Dead tumor centers, tough scar tissue, or tiny moving spots can yield unclear tissue cores. Aiming at the living outer edge using fused CT or ultrasound maps and checking needle position from two angles keeps accuracy high.
    \item \textbf{Next Steps for Unclear Results:} If a run or biopsy is inconclusive, your doctor may order a specialized CT angiogram, sample another viable spot, or discuss options with surgery and cancer specialists.
\end{itemize}

\section{Augmenting Methods}

\begin{itemize}
    \item \textbf{3D Cone-Beam CT and Scan Fusion:} Combines live 3D C-arm scans with earlier CT, MRI, or PET scans right in the procedure room to confirm tumor-feeding vessels and verify complete tumor treatment margins.
    \item \textbf{Tiny Inside-the-Vessel Cameras (IVUS/OCT):} Hair-thin ultrasound or light catheters measure exact vessel size, plaque buildup, and stent fit from the inside without extra X-rays or dye.
    \item \textbf{Before-and-After CT, MRI, and PET Scans:} Map anatomy before treatment and measure how well tumors or vessels responded afterward.
    \item \textbf{Pathology, Genetics, and Lab Tests:} Biopsy tissue tests, DNA profiling, infection cultures, and blood markers (like AFP, bilirubin, and kidney labs) guide targeted cancer, infection, and vessel care.
\end{itemize}

\section{Comparison with Alternative Modalities}

\begin{itemize}
    \item \textbf{Compared with Open Surgery:} IR uses tiny pinhole openings, causes less pain and shorter hospital stays, and often avoids general anesthesia. Open surgery remains necessary when large tumors must be removed whole or complex organs need rebuilding.
    \item \textbf{Compared with Endoscopy (ERCP/EUS):} Endoscopy is usually first-line for bleeding inside the stomach or bowel and standard bile-duct access, while IR shines when past surgery changed the anatomy, endoscopy could not stop bleeding, or deep outside-the-bowel access is needed.
    \item \textbf{Compared with Standard Diagnostic CT/MRI:} Regular CT and MRI safely screen and stage disease without any skin puncture, while IR steps in to collect definitive biopsy tissue and deliver targeted pinhole treatment in the same visit.
\end{itemize}

\section{Follow-up Steps}

\begin{itemize}
    \item \textbf{Recovery and Entry-Site Checks:} After vessel procedures, you will rest flat briefly while nurses check your vital signs, pulses, and entry spot for any swelling or bleeding. Many IR procedures are performed on an outpatient basis, allowing you to go home the same day once sedation wears off and the puncture site is stable. After chest biopsies or drains, a quick chest X-ray may check for a small air leak, and a skin check is scheduled if high X-ray dose thresholds (over 5 Gy) were reached.
    \item \textbf{Medicines, Comfort, and Drain Care:} You will go home with clear instructions on when to restart blood thinners, how to manage mild temporary ``post-embolization'' aches, low fever, or nausea, and how to gently flush and keep dressings clean if you have a temporary drain tube.
    \item \textbf{Lab and Biopsy Results:} Your care team follows up on tissue biopsy results or infection cultures to guide next treatments, and checks blood counts, kidney labs, or liver labs when appropriate.
    \item \textbf{Follow-Up Visits and Healing Scans:} Expect a follow-up visit and checkup scan---such as a CT or MRI 4--6 weeks after tumor ablation or embolization, an ultrasound after a stent or TIPS, or a quick contrast check before removing a drain tube or IVC filter.
\end{itemize}

\section{Key References}

\begin{itemize}
    \item Patel IJ, Rahim S, Davidson JC, et al. Society of Interventional Radiology Consensus Guidelines for the Periprocedural Management of Thrombotic and Bleeding Risk in Patients Undergoing Percutaneous Image-Guided Interventions. \textit{Journal of Vascular and Interventional Radiology}, 2019;30(8):1168--1184.
    \item Chehab MA, Thakor AS, Tulin-Silver S, et al. Adult and Pediatric Antibiotic Prophylaxis During Vascular and IR Procedures: A Society of Interventional Radiology Practice Parameter Update. \textit{Journal of Vascular and Interventional Radiology}, 2018;29(11):1483--1501.
    \item Miller DL, Balter S, Dixon RG, et al. Quality Improvement Guidelines for Recording Patient Radiation Dose in the Medical Record for Fluoroscopically Guided Procedures (SIR/CIRSE). \textit{Journal of Vascular and Interventional Radiology}, 2012;23(1):11--18.
    \item Kaufman JA, Lee MJ. \textit{Vascular and Interventional Radiology: The Requisites}, 3rd Edition. Elsevier, 2023.
    \item Stanford Health Care. \textit{Interventional Radiology}. Stanford Medicine. Available at: \url{https://stanfordhealthcare.org/medical-clinics/interventional-radiology.html}.
    \item Radiopaedia.org. \textit{Interventional radiology}. Radiopaedia. Available at: \url{https://radiopaedia.org/articles/interventional-radiology}.
\end{itemize}

\section{Further Reading}
\begin{itemize}
  \item Wikipedia contributors. \textit{Interventional radiology}. Wikipedia, The Free Encyclopedia. Available at: \url{https://en.wikipedia.org/wiki/Interventional_radiology}.
\end{itemize}
